\chapter{Diseño y Cálculos}
\label{chap:diseno-calculos}

El sistema se divide en cuatro etapas principales, cada una encargada de generar una de las señales requeridas.

\section{Etapa 1: Generación de Señal 1 (Cuadrada)}
Se requiere una onda cuadrada de $\SI{1.5}{\volt}$ a $\SI{-1}{\volt}$. 
La impedancia de entrada debe ser de $\SI{15}{\kilo\ohm}$.

% TODO: Agregar cálculos de la red de realimentación y offset para esta etapa.

\section{Etapa 2: Generación de Señal 2 (Triangular)}
Esta etapa recibe la señal cuadrada y realiza una integración para obtener una onda triangular. 
Rango: $\SI{4}{\volt}$ a $\SI{-1.5}{\volt}$.
Se utiliza una referencia de precisión con buffer para establecer el nivel de continua.

% TODO: Agregar cálculos de la constante de tiempo RC del integrador.

\section{Etapa 3: Generación de Señal 3 (Pulsos)}
La señal 3 es una onda de pulsos de $\SI{3}{\volt}$ a $\SI{0}{\volt}$. Se obtiene mediante un comparador.

\section{Etapa 4: Generación de Señal 4 (Picos)}
Para obtener los picos de $\SI{1}{\volt}$, se emplea un circuito derivador o una red RC de constante de tiempo muy pequeña seguida de un rectificador o limitador.

\section{Impedancia de Salida del Sistema}
La impedancia de salida se determina analizando la última etapa configurada...
